\subsection{\textbf{Sprint 1}}

Após a apresentação com as \textit{stakeholders}, reunimos o time para dividir as \textit{squads}. Para esta \textit{sprint}, planejamos doze tarefas de desenvolvimento, sendo seis referentes ao Épico 3, de questões, e seis distribuídas entre os demais épicos. Como o número de entregas era menor que o de integrantes, algumas delas acabaram sendo feitas em dupla, normalmente com um \ac{ages} I e um \ac{ages} II trabalhando juntos. O meu time ficou responsável pelas tarefas de autenticação e cadastro do usuário. A distribuição das histórias dentro de cada \textit{squad} ficou a cargo do \ac{ages} IV responsável por ela.

Um detalhe importante é que, antes mesmo do início da \textit{sprint}, um dos \ac{ages} III já havia adiantado a parte de autenticação no \textit{backend}, cobrindo tanto o login quanto o cadastro do usuário. Isso fez com que essa frente já chegasse praticamente concluída na minha \textit{squad}. Entreguei para o \ac{ages} II do time a tarefa de \textit{frontend} de cadastro, e para o \ac{ages} I o carrossel de primeiro acesso ao aplicativo. Como meu \ac{ages} III já tinha adiantado parte do \textit{backend}, optei por passar a ele uma atividade menos complexa, o que permitiria que ele apoiasse os \ac{ages} menores, revisasse os \ac{pr} e começasse a desenhar o diagrama da \ac{aws}, uma entrega que precisaria sair o quanto antes.

Logo após a primeira aula da \textit{sprint}, percebemos que o engajamento do time estava bastante alto. Já havia \ac{pr} abertos para a \textit{develop} nos primeiros dias. Com a capacidade sobrando, começamos a puxar histórias da \textit{Sprint} 3, já que boa parte das entregas da \textit{Sprint} 2 dependia diretamente de cadastro, login e questões ainda em andamento na \textit{Sprint} 1. Antecipar histórias da \textit{Sprint} 3 era mais seguro, pois elas tinham menos acoplamento com o que estava sendo construído naquele momento.

No meu time, no entanto, notei que as coisas andavam um pouco mais devagar, especialmente na tarefa do \ac{ages} II, que estava enfrentando dificuldades com o próprio Git. Conversei individualmente com cada integrante da \textit{squad} para entender se havia algo em que eu pudesse ajudar. Todos responderam que estava tudo certo e que as atividades estavam encaminhadas.

Na aula seguinte, conversei com a \ac{ages} I da \textit{squad} e ela me disse que a tarefa dela estava pronta, mas que tinha dúvida se deveria subir mesmo faltando uma rota. No fim, a rota dependia de uma tela de outra \textit{squad} que acabou atrasando, o que bloqueou parcialmente a entrega dela, ainda que a parte que cabia à \ac{ages} I já estivesse pronta. Já o \ac{ages} II tinha um \ac{pr} aberto com comentários dos \ac{ages} III pedindo ajustes, todos pontuais e resolvidos em aula. A única pendência real que ficou foi a história do \ac{ages} III do time, que acabou atrasando. Como era uma entrega importante, o componente de navegação do aplicativo, que destravaria várias atividades futuras, pedi que outro \ac{ages} III pegasse a tarefa antes do \textit{code freeze}, que havíamos estipulado para uma semana antes da apresentação. O Gabriel Saldanha, de outra \textit{squad}, puxou a história para si e conseguiu finalizá-la faltando três dias para a apresentação, já dentro do \textit{code freeze}.

Na aula anterior à apresentação, tentamos rodar um ensaio interno do que havia sido desenvolvido ao longo da \textit{sprint}, mas não tivemos sucesso. O repositório do \textit{mobile} estava com erro ao tentar emular o app no Expo. Mais uma vez, o Gabriel Saldanha se mostrou prestativo e ficou responsável por ajustar o problema.

Depois dessa aula, me reuni com os outros \ac{ages} IV para montar a apresentação e fazer um último teste no fluxo completo. Foi nesse momento que notamos uma quantidade considerável de \textit{bugs} que haviam passado para a \textit{develop}, alguns bastante escancarados. Em caráter de emergência, acionamos todos os \ac{ages} III para verificar o que conseguiriam corrigir até a apresentação. No fim, a maior parte dos \textit{bugs} foi resolvida, mas alguns persistiram. Anotamos como ponto de ação a criação de histórias específicas para resolução desses \textit{bugs} na \textit{sprint} seguinte.

Na apresentação, mostramos às \textit{stakeholders} tudo o que havia sido desenvolvido e reforçamos com elas o envio dos materiais que ainda estavam pendentes desde a \textit{Sprint} 0, os ícones do aplicativo e a base de questões. Elas ficaram bastante satisfeitas com o que viram e se comprometeram a enviar o material em seguida.

Apesar dos contratempos, conseguimos entregar mais atividades nesta \textit{sprint} do que havíamos planejado, o que reforçou a leitura de que a maior parte do time estava bem engajada com o projeto.

Após a apresentação, fizemos a retrospectiva da \textit{sprint}. Muitos pontos positivos foram levantados em relação à organização e à comunicação entre os times, mas também apareceram pontos de atenção. Como já havia sido comentado na \textit{Sprint} 0, algumas histórias não estavam totalmente alinhadas com o modelo de banco, o que acabou gerando dúvidas durante o desenvolvimento, e algumas regras de negócio não ficaram explícitas o suficiente nas \ac{us}. Criamos algumas pendências internas entre os \ac{ages} IV para serem resolvidas até a próxima aula. Ficou combinado que cada \ac{ages} IV revisaria as histórias da \textit{Sprint} 2 do seu time e solicitaria uma segunda revisão de outro \ac{ages} IV, para garantir que as \ac{us} chegassem ao time com o nível de clareza esperado.

Como \ac{ages} IV, tive dificuldade em acompanhar de perto o \ac{ages} III da minha \textit{squad} nos momentos finais da \textit{sprint}. Nos dias que antecederam o \textit{code freeze}, tentei contato algumas vezes para entender como estava o andamento da história e oferecer apoio, mas tive pouco retorno, o que atrasou a decisão de redistribuir a entrega e acabou se tornando um risco concreto para a \textit{sprint}. Outro ponto que só ficou claro ao longo dessas semanas foi a dificuldade de planejar um projeto inteiro no período curto que temos, algo que impacta tanto o detalhamento das histórias quanto a previsão de dependências entre as frentes.
